\documentclass{article}
\usepackage[utf8]{inputenc}
\usepackage{a4wide}

\begin{document}

\section*{Nugget}

ร้านฟาสต์ฟู้ดแห่งหนึ่งขายนักเก็ตเป็นกล่อง มีกล่องนักเก็ตอยู่ 3 ขนาด คือ เล็ก, กลาง, และใหญ่

ใส่นักเก็ตจำนวน



6, 9, และ 20 ชิ้นตามลำดับ



เลขนักเก็ต คือจำนวนเต็มบวกที่เกิดจากผลรวมของจำนวนนักเก็ตในกล่องขนาดต่างๆ เช่น เลข 6

เป็นเลขนักเก็ต



เพราะเป็นจำนวนนักเก็ตในกล่องเล็ก, เลข 12

เป็นเลขนักเก็ตเพราะเกิดจากการรวมกันของจำนวนนักเก็ตในกล่อง



เล็กสองกล่อง, เลข 15

เป็นเลขนักเก็ตเพราะเกิดจากการรวมกันของจำนวนนักเก็ตในกล่องเล็กหนึ่งกล่องและกล่อง



กลางหนึ่งกล่อง เป็นต้น เลข 4 และ 10

ไม่เป็นเลขนักเก็ตเพราะเลขดังกล่าวไม่สามารถเกิดจากการรวมกันของจำนวน



นักเก็ตในกล่องขนาดใดๆ ได้



\subsection*{Input}
{รับค่า n ที่เป็นจำนวนเต็ม จาก standard input โดยที่ 1 ≤ n ≤ 100}\\



\subsection*{Output}
มีหลายบรรทัด พิมพ์เลขนักเก็ตที่น้อยกว่าหรือเท่ากับ n โดยเรียงค่าจากน้อยไปหามาก

พิมพ์บรรทัดละหนึ่งตัวเลข



ถ้าไม่มีเลขนักเก็ตที่น้อยกว่าหรือเท่ากับ n ให้พิมพ์คำว่า "no"



ถ้าพิมเกิน 100 ให้พิมว่าInput exceeds the maximum limit of 100



\end{document}
